% ----------------------------------------------------
\begin{frame}
  \titlepage
\end{frame}
% ----------------------------------------------------
\note{}

% ----------------------------------------------------
\begin{frame}
\frametitle{Today's goals}
\centering

\begin{itemize}[<+->]
\item Understand why computing practices matter for research
\item Learn principles for organizing a research project
\item Integrate R output directly into \LaTeX{} documents
\item Write better R code: functions, checks, and style
\item Understand the role of plain text and command line tools
\item Deepen your knowledge of version control with Git
\end{itemize}

\end{frame}
% ----------------------------------------------------
\note{This is a practical session focused on the computing side of doing empirical research. The skills covered today are not statistical, but they are essential for doing good empirical work. Students who have struggled with organizing their assignments or tracking down bugs in their code will find this session especially useful. The goal is not to master all of these tools in one session, but to establish good habits that will pay off for the rest of the course and beyond.}

% ====================================================
\section{Introduction}
% ====================================================

% ----------------------------------------------------
\begin{frame}
\frametitle{Code is a scientific artifact}
\centering

\begin{itemize}[<+->]
\item Social scientists often treat computing as a necessary evil
  \begin{itemize}
  \item Something to get through to reach results
  \end{itemize}
\vspace{8pt}
\item But code is part of your research
  \begin{itemize}
  \item It will be \textbf{read} (by collaborators, reviewers, future you)
  \item It will \textbf{break} (after software updates, on another machine)
  \item It needs to be \textbf{maintained} (for revisions, new data)
  \end{itemize}
\vspace{8pt}
\item The standards of rigor you apply to your theory and research design should extend to your computing
\end{itemize}

\end{frame}
% ----------------------------------------------------
\note{This is the key mindset shift for the whole lecture. Most students in the social sciences have been trained to think carefully about research design and causal inference, but treat code as a throwaway scratchpad --- something they write once, get results from, and never look at again. The reality is that code is a scientific artifact: it embodies your analytical choices, it is the thing that actually produces your findings, and it will be read by others (reviewers, co-authors, or your future self trying to revise the paper). If the code is messy, undocumented, or unreproducible, the research is fragile regardless of how good the identification strategy is.}

% ----------------------------------------------------
\begin{frame}
\frametitle{Why computing practices matter}
\centering

\begin{itemize}[<+->]
\item \textbf{Reproducibility}: can you (or someone else) re-run your analysis?
\vspace{8pt}
\item \textbf{Avoiding errors}: manual steps introduce mistakes
  \begin{itemize}
  \item Copy-pasting results into Word
  \item Renaming files by hand
  \item Running scripts in the wrong order
  \end{itemize}
\vspace{8pt}
\item \textbf{Efficiency}: time invested in workflow saves time later
\vspace{8pt}
\item On errors: \href{https://journals.sagepub.com/doi/10.1177/20531680221126454}{Bisbee et al.\ (2022)}
\end{itemize}

\end{frame}
% ----------------------------------------------------
\note{Start by motivating why we are spending a whole session on computing practices. The reproducibility crisis in social science is partly a computing problem: researchers cannot reproduce their own results because they lost track of which script produced which output, or because they made manual changes that were never recorded. The link at the bottom is a paper documenting how common coding errors are in published research, including errors introduced by tools like stargazer that silently reorder variables. Emphasize that these are not hypothetical problems --- they happen to experienced researchers regularly.}

% ----------------------------------------------------
\begin{frame}
\centering
\vspace{30pt}
{\large You come back to a project after 6 months.\\\vspace{15pt}You need to update one figure.\\\vspace{15pt}How long does it take you?}
\end{frame}
% ----------------------------------------------------
\note{Discussion prompt. Ask students to think about their own experience. Most will admit that going back to an old project is painful: they do not remember which script produces which output, what the file names mean, or what order to run things in. The goal of this session is to make the answer to this question ``a few minutes'' instead of ``a whole day.'' A well-organized project with a Makefile and clear folder structure can be re-run with a single command. A poorly organized project requires archaeology.}

% ====================================================
\section{Project Organization}
% ====================================================

% ----------------------------------------------------
\begin{frame}
\frametitle{Three principles}
\centering

\begin{itemize}[<+->]
\item \textbf{Separation of concerns}
  \begin{itemize}
  \item Code does analysis. Documents present it.
  \item Output files (tables, figures) are the interface between the two
  \end{itemize}
\vspace{8pt}
\item \textbf{The pipeline}
  \begin{itemize}
  \item Raw data $\rightarrow$ cleaned data $\rightarrow$ analysis $\rightarrow$ results $\rightarrow$ document
  \item Each stage has clear inputs and outputs
  \item No stage depends on a later stage
  \end{itemize}
\vspace{8pt}
\item \textbf{Raw data is sacred}
  \begin{itemize}
  \item Never overwrite it, never modify it by hand
  \item Code transforms raw data into everything else
  \end{itemize}
\end{itemize}

\end{frame}
% ----------------------------------------------------
\note{These three principles are the foundation of good project organization. Separation of concerns means that your analysis scripts do not contain LaTeX code, and your documents do not contain R code --- the only connection between them is the output files (a .tex table, a .pdf figure). The pipeline principle means your project has a clear direction: data flows from raw to cleaned to analyzed to presented, and no step depends on a later step. This makes the project debuggable --- if a table looks wrong, you know to check the analysis script, not the document. Raw data immutability is the most important rule: if you modify your raw data by hand, you cannot reproduce your results. All transformations must happen in code, so they are documented and repeatable.}

% ----------------------------------------------------
\begin{frame}[fragile]
\frametitle{Organizing a project: one example}
\centering

\vspace{5pt}

{\footnotesize\ttfamily
\begin{tabular}{@{}l@{}}
\textcolor{accent}{\textbf{workflow\_example/}} \\
\hspace{1em}Makefile \\
\hspace{1em}.gitignore \\
\hspace{1em}\textcolor{accent}{\textbf{create\_data/}} \\
\hspace{2em}data.R \\
\hspace{2em}\textcolor{asher}{output/}\,data.csv \\
\hspace{1em}\textcolor{accent}{\textbf{analyses/}} \\
\hspace{2em}analyze.R \\
\hspace{2em}\textcolor{asher}{output/}\,table\_models.tex \\
\hspace{1em}\textcolor{accent}{\textbf{plots/}} \\
\hspace{2em}plots.R \\
\hspace{2em}\textcolor{asher}{output/}\,scatter.pdf \\
\hspace{1em}\textcolor{accent}{\textbf{document/}} \\
\hspace{2em}doc.tex \\
\end{tabular}
}

\vspace{10pt}

\begin{itemize}
\item Full example: \href{https://github.com/franvillamil/workflow_example}{\texttt{github.com/franvillamil/workflow\_example}}
\end{itemize}

\end{frame}
% ----------------------------------------------------
\note{This is one way to implement those principles. Each folder is a task: create\_data generates the dataset, analyses runs models and produces a LaTeX table, plots creates figures, and document compiles the final paper. Each task folder has a script and an output subfolder. The Makefile ties everything together (we will see it shortly). The .gitignore tells Git which files to ignore. The key insight: the document folder does not contain any R code, and the analysis folders do not contain any LaTeX. Each piece does one thing. There are other valid ways to organize a project --- the details matter less than the principles: clear separation, a pipeline that flows in one direction, and raw data that is never touched.}

% ----------------------------------------------------
\begin{frame}
\frametitle{Practical rules}
\centering

\begin{itemize}[<+->]
\item \textbf{Always use relative paths}
  \begin{itemize}
  \item \red{\texttt{read.csv("/Users/fran/proj/data/file.csv")}} --- breaks on any other machine
  \item \textcolor{accent}{\texttt{read.csv("data/file.csv")}} --- works anywhere
  \end{itemize}
\vspace{8pt}
\item \textbf{Never one giant script}
  \begin{itemize}
  \item Break the project into discrete steps
  \item Each script has clear inputs and outputs
  \end{itemize}
\vspace{8pt}
\item \textbf{Generated files are disposable}
  \begin{itemize}
  \item If a file can be reproduced by running a script, it is not a source file
  \item Delete all output, re-run, and get the same results
  \end{itemize}
\end{itemize}

\end{frame}
% ----------------------------------------------------
\note{Relative paths are the most common reason code fails on another machine. If your script says \texttt{read.csv("/Users/yourname/Desktop/project/data.csv")}, it will not work on anyone else's computer. Use relative paths from the project root: \texttt{read.csv("data/file.csv")}. The ``never one giant script'' rule is about modularity: a 500-line script that loads data, cleans it, runs models, and makes plots is hard to debug and hard to maintain. Break it into steps. The ``generated files are disposable'' rule is a test: can you delete everything in your output folders, re-run your scripts, and get exactly the same results? If yes, your project is reproducible. If not, something is missing from your pipeline.}

% ----------------------------------------------------
\begin{frame}
\frametitle{File naming conventions}
\centering

\vspace{10pt}

{\footnotesize
\begin{tabular}{ll}
\hline
\textbf{\red{Bad}} & \textbf{\textcolor{accent}{Good}} \\
\hline
\texttt{Final Data.csv} & \texttt{data\_cleaned.csv} \\
\texttt{Datos educaci\'{o}n.csv} & \texttt{datos\_educacion.csv} \\
\texttt{analysis.R} (which one?) & \texttt{01\_clean\_data.R} \\
\texttt{figure1final2.pdf} & \texttt{fig\_scatter\_income.pdf} \\
\texttt{My Thesis Draft (3).docx} & \texttt{thesis\_draft.tex} \\
\hline
\end{tabular}
}

\vspace{15pt}

\begin{itemize}[<+->]
\item \textbf{No spaces} in file names (use underscores or hyphens)
\item \textbf{No special characters} (accents, symbols)
\item Use \textbf{numbered prefixes} for scripts that run in order
\item Use \textbf{descriptive names}: what is in the file, not when you made it
\end{itemize}

\end{frame}
% ----------------------------------------------------
\note{File naming seems trivial but causes real problems. Spaces in file names break command-line tools and require quoting in scripts. Special characters cause encoding issues across operating systems. Numbered prefixes (01\_, 02\_) make the execution order obvious. Descriptive names mean you can find the right file without opening it. The ``Final Data (3).csv'' pattern is a sign of a disorganized project. If you need version control, use Git, not file names. The goal is that anyone (including your future self) can look at the file listing and understand what each file does.}

% ----------------------------------------------------
\begin{frame}
\frametitle{Integrating R and \LaTeX}
\centering

\begin{itemize}[<+->]
\item R produces output files (tables, figures)
\item \LaTeX{} \texttt{\textbackslash input\{\}} or \texttt{\textbackslash includegraphics\{\}} reads them
\vspace{8pt}
\item Example from \texttt{workflow\_example}:
  \begin{itemize}
  \item \texttt{analyze.R} $\rightarrow$ \texttt{analyses/output/table\_models.tex}
  \item \texttt{plots.R} $\rightarrow$ \texttt{plots/output/scatter.pdf}
  \item \texttt{doc.tex}: \texttt{\textbackslash input\{../analyses/output/table\_models\}}
  \item \texttt{doc.tex}: \texttt{\textbackslash includegraphics\{../plots/output/scatter\}}
  \end{itemize}
\vspace{8pt}
\item \textbf{No copy-pasting}: update the analysis, recompile the document
\item No local \LaTeX? Use \textbf{Overleaf} --- same \texttt{\textbackslash input\{\}} workflow
\end{itemize}

\end{frame}
% ----------------------------------------------------
\note{This is one of the biggest practical payoffs of good project organization. When your LaTeX document reads tables and figures directly from the output folders, you never have to manually update them. Change a model specification? Re-run the R script, recompile the document, and everything updates automatically. This eliminates a major source of errors: the table in the paper not matching the actual analysis because you forgot to update it after a change. For students who do not have LaTeX installed locally: Overleaf works the same way. You can upload the generated .tex table files and .pdf figures to your Overleaf project and use \texttt{\textbackslash input\{\}} and \texttt{\textbackslash includegraphics\{\}} just the same. Overleaf also supports Git, so you can clone your Overleaf project locally and copy output files into it programmatically.}

% ----------------------------------------------------
\begin{frame}
\frametitle{Automating the pipeline: Makefiles}
\centering

\vspace{5pt}

{\scriptsize
\begin{itemize}
\item[] \texttt{all: create\_data/output/data.csv analyses/output/table\_models.tex \textbackslash}
\item[] \texttt{\ \ \ \ plots/output/scatter.pdf}
\item[] \texttt{create\_data/output/data.csv: create\_data/data.R}
\item[] \texttt{\ \ \ \ Rscript --no-save create\_data/data.R}
\item[] \texttt{analyses/output/table\_models.tex: analyses/analyze.R \textbackslash}
\item[] \texttt{\ \ \ \ \ \ \ \ create\_data/output/data.csv}
\item[] \texttt{\ \ \ \ Rscript --no-save analyses/analyze.R}
\item[] \texttt{plots/output/scatter.pdf: plots/plots.R \textbackslash}
\item[] \texttt{\ \ \ \ \ \ \ \ create\_data/output/data.csv}
\item[] \texttt{\ \ \ \ Rscript --no-save plots/plots.R}
\end{itemize}
}

\vspace{5pt}
{\footnotesize \textbf{Note:} Makefiles require \textbf{tabs} (not spaces) for indentation}

\end{frame}
% ----------------------------------------------------
\note{Walk through the Makefile line by line. The first rule (all) lists all the final outputs. Each subsequent rule has a target (the output file), dependencies (the script and any input files), and a recipe (the command to run). The key feature: Make only re-runs a step if its dependencies have changed. If you modify plots.R but not analyze.R, only the plots step re-runs. This saves time on large projects and ensures consistency. To run everything: type \texttt{make} in the terminal. To run just one step: \texttt{make plots/output/scatter.pdf}. Makefiles require tabs (not spaces) for indentation --- this is a common source of errors. Reference: \texttt{makefiletutorial.com}.}

% ----------------------------------------------------
\begin{frame}
\frametitle{Makefiles: the logic}
\centering

\vspace{10pt}

\begin{tikzpicture}[scale=0.78, every node/.style={transform shape},
  box/.style={draw, rounded corners, minimum width=2.4cm, minimum height=0.8cm, align=center, font=\footnotesize\ttfamily},
  arrow/.style={->, thick, accent}
]
  % Nodes
  \node[box, fill=accent!10] (data) at (0,0) {data.R};
  \node[box, fill=accent!10] (csv) at (3.5,0) {data.csv};
  \node[box, fill=accent!10] (analyze) at (7,1.2) {analyze.R};
  \node[box, fill=accent!10] (plots) at (7,-1.2) {plots.R};
  \node[box, fill=accent2!10] (table) at (10.5,1.2) {table.tex};
  \node[box, fill=accent2!10] (scatter) at (10.5,-1.2) {scatter.pdf};
  % Arrows
  \draw[arrow] (data) -- (csv);
  \draw[arrow] (csv) -- (analyze);
  \draw[arrow] (csv) -- (plots);
  \draw[arrow] (analyze) -- (table);
  \draw[arrow] (plots) -- (scatter);
  % Labels
  \node[above, font=\scriptsize, asher] at (1.75,0.15) {produces};
  \node[above, font=\scriptsize, asher] at (8.75,1.35) {produces};
  \node[above, font=\scriptsize, asher] at (8.75,-1.05) {produces};
\end{tikzpicture}

\vspace{15pt}

\begin{itemize}
\item \texttt{make} only re-runs what has \textbf{changed}
\item The dependency graph ensures correct ordering
\end{itemize}

\end{frame}
% ----------------------------------------------------
\note{This diagram shows the dependency graph that Make constructs from the Makefile. Each arrow represents a dependency: the analysis script depends on the data, and the table depends on the analysis script. If you change data.R, Make knows it needs to re-run data.R, then analyze.R, then plots.R (because they all depend on data.csv). If you only change plots.R, only the plots step re-runs. This is the power of Make: it tracks dependencies automatically. For small projects, running everything manually is fine. For larger projects (a dissertation with 10 scripts), Make saves hours and prevents errors from running things in the wrong order.}

% ====================================================
\section{Writing Better Code}
% ====================================================

% ----------------------------------------------------
\begin{frame}
\frametitle{DRY: Don't Repeat Yourself}
\centering

\begin{itemize}[<+->]
\item If you write the same code \textbf{twice}, write a \textbf{function}
\vspace{8pt}
\item Why?
  \begin{itemize}
  \item Fix a bug in one place, not ten
  \item Easier to test and debug
  \item Shorter, cleaner scripts
  \end{itemize}
\vspace{8pt}
\item Common candidates for functions:
  \begin{itemize}
  \item Cleaning steps applied to multiple datasets
  \item Running the same model with different variables
  \item Producing plots with consistent formatting
  \end{itemize}
\end{itemize}

\end{frame}
% ----------------------------------------------------
\note{The DRY principle is the single most important coding habit. Whenever you find yourself copying and pasting code and changing one or two things, you should write a function instead. The reason is not just aesthetics: duplicated code means duplicated bugs. If you find an error in the cleaning step for one dataset, you have to remember to fix it in all the copies. With a function, you fix it once. Functions also make your code self-documenting: \texttt{clean\_survey(df)} is more readable than 20 lines of dplyr operations. Encourage students to start small: even a function that just wraps a few lines is better than copy-pasting.}

% ----------------------------------------------------
\begin{frame}
\frametitle{Define constants at the top}
\centering

\vspace{10pt}

{\footnotesize
\begin{tabular}{ll}
\hline
\textbf{\red{Bad}} & \textbf{\textcolor{accent}{Good}} \\
\hline
\texttt{df = df[df\$year >= 2000, ]} & \texttt{start\_year = 2000} \\
\texttt{...} & \texttt{...} \\
\texttt{df2 = df2[df2\$year >= 2000, ]} & \texttt{df = df[df\$year >= start\_year, ]} \\
 & \texttt{df2 = df2[df2\$year >= start\_year, ]} \\
\hline
\end{tabular}
}

\vspace{15pt}

\begin{itemize}[<+->]
\item Put parameters and thresholds at the \textbf{top of the script}
\item Change once, applies everywhere
\vspace{8pt}
\item From the example project:
  \begin{itemize}
  \item \texttt{n\_obs = 1000}
  \item Used throughout --- change it once to re-run with different sample size
  \end{itemize}
\end{itemize}

\end{frame}
% ----------------------------------------------------
\note{This is a simple but powerful habit. Any number or string that appears more than once in your script should be defined as a constant at the top. Year cutoffs, sample size, file paths, variable names for analysis --- all of these should be constants. This way, when you need to change a parameter (e.g., switch from 2000 to 1990), you change it in one place. The alternative --- searching through a 200-line script for every occurrence of ``2000'' --- is error-prone and tedious. In the workflow\_example project, \texttt{n\_obs = 1000} is defined once and used in the data generation step. Changing it to 5000 requires editing one line.}

% ----------------------------------------------------
\begin{frame}
\frametitle{Write checks and assertions}
\centering

\vspace{5pt}

{\scriptsize
\begin{itemize}\setlength{\itemsep}{0pt}
\item[] \texttt{\# After merging two datasets}
\item[] \texttt{merged = merge(df1, df2, by = "id")}
\item[] \texttt{if(nrow(merged) != nrow(df1)) \{}
\item[] \texttt{\ \ stop("Merge changed number of rows!")}
\item[] \texttt{\}}
\item[]
\item[] \texttt{\# Check for duplicates}
\item[] \texttt{if(any(duplicated(df\$id))) \{}
\item[] \texttt{\ \ stop("Duplicate IDs found!")}
\item[] \texttt{\}}
\item[]
\item[] \texttt{\# Sanity check on values}
\item[] \texttt{if(any(df\$age < 0 | df\$age > 120, na.rm = TRUE)) \{}
\item[] \texttt{\ \ warning("Suspicious age values detected")}
\item[] \texttt{\}}
\end{itemize}
}

\end{frame}
% ----------------------------------------------------
\note{Assertions are the best defense against silent errors. The idea: after every operation that could go wrong, add a check that verifies the result is what you expect. Merges are the classic example: a merge that unexpectedly creates duplicates can silently inflate your sample size and bias your results. The \texttt{stop()} function halts execution with an error message --- use it when something is definitely wrong. The \texttt{warning()} function prints a warning but continues --- use it when something is suspicious but might be okay. These checks cost nothing to write and can save you from publishing wrong results. Make it a habit to add a check after every merge, filter, or reshape operation.}

% ----------------------------------------------------
\begin{frame}
\frametitle{Print diagnostics as you go}
\centering

\begin{itemize}[<+->]
\item Sprinkle \texttt{print()} and \texttt{cat()} throughout your scripts
\vspace{8pt}
\item Examples:
  \begin{itemize}
  \item \texttt{cat("Rows after cleaning:", nrow(df), "\textbackslash n")}
  \item \texttt{print(table(df\$treatment))}
  \item \texttt{cat("Missing values:", sum(is.na(df\$outcome)), "\textbackslash n")}
  \end{itemize}
\vspace{8pt}
\item When the script runs, you get a \textbf{log} of what happened
\item If something goes wrong later, the log tells you where
\vspace{8pt}
\item This is especially valuable when running scripts via \texttt{Rscript} or \texttt{make}
\end{itemize}

\end{frame}
% ----------------------------------------------------
\note{Print statements are the simplest form of logging. When you run a script interactively in RStudio, you can inspect objects as you go. But when you run scripts via the command line (which is what Make does), you do not see intermediate results unless you print them. Adding print statements that show sample sizes, value distributions, and missing data counts creates a log that you can review to make sure everything went as expected. This is especially important for long-running scripts that process multiple datasets. If the final output looks wrong, the print log helps you narrow down where the problem occurred.}

% ----------------------------------------------------
\begin{frame}
\frametitle{Seeds and package versions}
\centering

\begin{itemize}[<+->]
\item Any code involving randomness must set a \textbf{seed}:
  \begin{itemize}
  \item Simulation, bootstrapping, random sampling
  \item \texttt{set.seed(12345)} at the top of the script
  \item Same seed $=$ same results every time
  \end{itemize}
\vspace{8pt}
\item Results can change across \textbf{package versions}
  \begin{itemize}
  \item Document which version of R and packages you used
  \item Check your R session info with \texttt{sessionInfo()}
  \item Options: \texttt{renv} (R) locks exact package versions
  \end{itemize}
\vspace{8pt}
\item Not paranoia --- this has caused real problems in published work
\end{itemize}

\end{frame}
% ----------------------------------------------------
\note{This sounds like a minor technical detail, but forgetting to set a seed is one of the most common reproducibility failures. If your analysis involves any randomness --- bootstrapped standard errors, simulations, cross-validation, random sampling --- the results will differ every time you run the script unless you set a seed. The seed should go at the top of the script so it is obvious. Package versions are a subtler problem: a function that behaved one way in version 1.2 might behave differently in version 2.0. The \texttt{renv} package in R creates a lockfile that records exact package versions, so someone else can install the same versions and reproduce your results. You do not need to use \texttt{renv} for every project, but you should at least note your R version and key package versions somewhere (e.g., in the README).}

% ----------------------------------------------------
\begin{frame}
\frametitle{Code style matters}
\centering

\begin{itemize}[<+->]
\item Use \textbf{meaningful variable names}:
  \begin{itemize}
  \item \texttt{pop\_2020} not \texttt{x1}, \texttt{income\_log} not \texttt{v}
  \end{itemize}
\vspace{8pt}
\item Be \textbf{consistent}:
  \begin{itemize}
  \item Pick \texttt{snake\_case} and stick with it
  \item Consistent indentation (2 spaces in R)
  \end{itemize}
\vspace{8pt}
\item Write \textbf{comments} for non-obvious code
  \begin{itemize}
  \item Why, not what: \texttt{\# Drop obs before 1990 (pre-reform)}
  \item Use section dividers: \texttt{\# ============}
  \end{itemize}
\vspace{8pt}
\item Your code is read more often than it is written
\end{itemize}

\end{frame}
% ----------------------------------------------------
\note{Code style is about making your code readable --- for others and for your future self. Meaningful variable names eliminate the need for comments explaining what \texttt{x1} is. Consistent formatting (indentation, spacing, naming conventions) makes the code scannable. Comments should explain the reasoning behind non-obvious decisions, not narrate what the code does line by line: ``Drop observations before 1990 per Acemoglu (2001)'' is useful; ``filter rows'' is not. The section divider pattern (a line of equals signs or dashes) helps break long scripts into navigable chunks. The R community has adopted snake\_case as the standard; avoid dots in names (data.frame style) and camelCase (Java style). Reference: Hadley Wickham's R Style Guide.}

% ----------------------------------------------------
\begin{frame}
\frametitle{What NOT to do}
\centering

\begin{itemize}[<+->]
\item Overwriting raw data with your code
\vspace{8pt}
\item Results that require clicking through a GUI
  \begin{itemize}
  \item SPSS menus, Excel transformations, manual edits
  \end{itemize}
\vspace{8pt}
\item ``It works on my machine''
  \begin{itemize}
  \item Almost always an absolute path problem
  \end{itemize}
\vspace{8pt}
\item ``I fixed a few values in Excel''
  \begin{itemize}
  \item Undocumented, unreproducible, invisible to collaborators
  \end{itemize}
\vspace{8pt}
\item Scripts that must be run in a specific undocumented order
\end{itemize}

\end{frame}
% ----------------------------------------------------
\note{Sometimes the most memorable part of a lecture is the horror stories. Each of these is a real pattern that leads to broken research. Overwriting raw data means you cannot go back to the original if something went wrong. GUI-dependent results (clicking through SPSS or Excel) cannot be reproduced because there is no record of what you did. Absolute paths are the number one reason code fails on another machine. Manual data edits are invisible: your co-author does not know you changed three values in Excel, and neither will the reviewer or your future self. Undocumented script ordering means the project only works if you happen to remember to run things in the right sequence --- and you will not remember in six months. Every one of these problems is solved by the practices covered in this lecture.}

% ====================================================
\section{Plain Text and Tools}
% ====================================================

% ----------------------------------------------------
\begin{frame}
\frametitle{Why plain text?}
\centering

\begin{itemize}[<+->]
\item Plain text files: \texttt{.R}, \texttt{.tex}, \texttt{.csv}, \texttt{.md}, \texttt{.py}
\item Binary files: \texttt{.docx}, \texttt{.xlsx}, \texttt{.dta}, \texttt{.pdf}
\vspace{8pt}
\item Plain text advantages:
  \begin{itemize}
  \item \textbf{Portable}: works on any computer, any OS, forever
  \item \textbf{Version control}: Git tracks changes line by line
  \item \textbf{Scriptable}: can be processed by other tools
  \item \textbf{No vendor lock-in}: does not require specific software
  \end{itemize}
\vspace{8pt}
\item You are already using plain text: R scripts, \LaTeX{} files
\item The goal: use it for \textbf{as much as possible}
\end{itemize}

\end{frame}
% ----------------------------------------------------
\note{Plain text is the foundation of a reproducible workflow. A .tex file from 1990 still compiles today; a Word 95 .doc file may not open in modern Word. Git can show you exactly which line changed between two versions of a .R file; for a .docx file, it can only tell you the file changed. CSV files can be read by any programming language; .xlsx files require special libraries. The point is not that binary formats are bad --- sometimes you need them (e.g., images, compiled PDFs). The point is that your source materials (code, text, data when possible) should be in plain text so they are portable, versionable, and scriptable. Reference: Kieran Healy's \textit{The Plain Person's Guide to Plain Text Social Science} at \texttt{plain-text.co}.}

% ----------------------------------------------------
\begin{frame}
\frametitle{Code editors}
\centering

\begin{itemize}[<+->]
\item RStudio is fine for R, but consider a \textbf{general-purpose editor}
\vspace{8pt}
\item Options:
  \begin{itemize}
  \item \textbf{VS Code}: free, huge ecosystem of extensions, very popular
  \item \textbf{Positron}: new IDE by Posit (RStudio makers), for R and Python
  \item \textbf{Sublime Text}: fast, lightweight, highly customizable
  \end{itemize}
\vspace{8pt}
\item Why use a general editor?
  \begin{itemize}
  \item Edit R, \LaTeX{}, Python, Makefiles in the same tool
  \item Better project navigation and search
  \item Customizable snippets and shortcuts
  \end{itemize}
\end{itemize}

\end{frame}
% ----------------------------------------------------
\note{Most students in the course use RStudio, which is a perfectly good tool for R. But as they start working with LaTeX, Makefiles, and possibly Python, having a general-purpose editor becomes valuable. VS Code is currently the most popular editor and has excellent R support through extensions. Positron is a new option from the same company that makes RStudio, designed for both R and Python. Sublime Text is lightweight and fast. The key message is not ``stop using RStudio'' but ``be aware that other tools exist and might serve you better as your workflow becomes more complex.'' Emphasize that the choice of editor is personal --- what matters is being comfortable and productive with your tool.}

% ----------------------------------------------------
\begin{frame}
\frametitle{The command line: basics}
\centering

\begin{itemize}[<+->]
\item The command line is how you talk directly to your computer
\vspace{8pt}
\item Essential commands:
  \begin{itemize}
  \item \texttt{cd} --- change directory
  \item \texttt{ls} --- list files
  \item \texttt{pwd} --- print working directory
  \item \texttt{mkdir} --- create a directory
  \item \texttt{Rscript file.R} --- run an R script
  \item \texttt{make} --- run a Makefile
  \end{itemize}
\vspace{8pt}
\item You \textbf{already use it} for Git (\texttt{git add}, \texttt{git commit}, \texttt{git push})
\item Terminal on Mac/Linux, Git Bash or WSL on Windows
\end{itemize}

\end{frame}
% ----------------------------------------------------
\note{Many students are intimidated by the command line, but they have already been using it for Git since Session 1. The goal here is not to turn them into command-line experts, but to make them comfortable with basic navigation and script execution. Being able to \texttt{cd} into a project folder and type \texttt{make} is all they need for the workflow we are teaching. The six commands listed here cover 90\% of what a social science researcher needs. For those who want to go deeper, point them to MIT's Missing Semester course (\texttt{missing.csail.mit.edu}) and Software Carpentry's Unix Shell lesson.}

% ====================================================
\section{Version Control with Git}
% ====================================================

% ----------------------------------------------------
\begin{frame}
\frametitle{Git: recap from Session 1}
\centering

\begin{itemize}[<+->]
\item You have been using Git since the first session
\vspace{8pt}
\item The basic workflow:
  \begin{itemize}
  \item \texttt{git add file.R} --- stage changes
  \item \texttt{git commit -m "message"} --- save a snapshot
  \item \texttt{git push} --- upload to GitHub
  \item \texttt{git pull} --- download from GitHub
  \end{itemize}
\vspace{8pt}
\item Today: going deeper into \textbf{good practices}
\end{itemize}

\end{frame}
% ----------------------------------------------------
\note{Quick recap of what students already know. By this point they should be comfortable with the basic add-commit-push cycle. Today we go deeper into the practices that make Git actually useful rather than just a backup tool. The three topics: .gitignore (what not to track), commit messages (how to write useful history), and GitHub for sharing replication materials.}

% ----------------------------------------------------
\begin{frame}
\frametitle{Git: commits as snapshots}
\centering

\vspace{5pt}

\begin{tikzpicture}[
  commit/.style={circle, draw=accent, fill=accent!20, thick, minimum size=0.7cm, font=\scriptsize\ttfamily},
  arrow/.style={->, thick, accent!70!black},
  label/.style={font=\scriptsize, text=asher, align=center}
]
  \node[commit] (c1) at (0,0) {c1};
  \node[commit] (c2) at (2.2,0) {c2};
  \node[commit] (c3) at (4.4,0) {c3};
  \node[commit] (c4) at (6.6,0) {c4};
  \node[commit, fill=accent2!20, draw=accent2] (c5) at (8.8,0) {c5};

  \draw[arrow] (c1) -- (c2);
  \draw[arrow] (c2) -- (c3);
  \draw[arrow] (c3) -- (c4);
  \draw[arrow] (c4) -- (c5);

  \node[label] at (0,1.0) {Initial\\ commit};
  \node[label] at (2.2,1.0) {Add data\\ cleaning};
  \node[label] at (4.4,1.0) {Run main\\ models};
  \node[label] at (6.6,1.0) {Fix merge\\ bug};
  \node[label] at (8.8,1.0) {Add robust\\ SE};

  \node[font=\scriptsize\bfseries, text=accent2] at (8.8,-0.9) {HEAD};
\end{tikzpicture}

\vspace{15pt}

\begin{itemize}
\item Each \textbf{commit} is a snapshot of your project at a point in time
\item The \textbf{history} is a chain of commits, each with a message
\item \texttt{HEAD} is a pointer to your current position
\end{itemize}

\end{frame}
% ----------------------------------------------------
\note{Visual intuition for what Git actually does. Each commit is a complete snapshot of the project state when you ran \texttt{git commit}. The chain of commits is your history. HEAD is a pointer that tells you where you are in that history --- normally at the tip of your current branch. When you run \texttt{git log}, you see this chain going backwards from HEAD. The key insight: Git does not store diffs between versions, it stores full snapshots (efficiently, via hashing). This is why commits are cheap and you should make them often.}

% ----------------------------------------------------
\begin{frame}
\frametitle{Branches: parallel lines of work}
\centering

\vspace{5pt}

\begin{tikzpicture}[
  commit/.style={circle, draw=accent, fill=accent!20, thick, minimum size=0.55cm},
  fcommit/.style={circle, draw=accent2, fill=accent2!20, thick, minimum size=0.55cm},
  arrow/.style={->, thick, accent!70!black},
  farrow/.style={->, thick, accent2!70!black},
  label/.style={font=\scriptsize\bfseries, align=center}
]
  \node[commit] (m1) at (0.8,0) {};
  \node[commit] (m2) at (2.5,0) {};
  \node[commit] (m3) at (4.2,0) {};
  \node[commit] (m4) at (7.6,0) {};
  \node[commit] (m5) at (9.3,0) {};

  \node[fcommit] (f1) at (5.9,-1.3) {};
  \node[fcommit] (f2) at (7.6,-1.3) {};

  \draw[arrow] (m1) -- (m2);
  \draw[arrow] (m2) -- (m3);
  \draw[arrow, dashed] (m3) -- (m4);
  \draw[arrow] (m4) -- (m5);

  \draw[farrow] (m3) -- (f1);
  \draw[farrow] (f1) -- (f2);
  \draw[farrow] (f2) -- (m4);

  \node[label, text=accent] at (-0.1,0) {main};
  \node[label, text=accent2] at (4.2,-1.3) {new branch};
  \node[font=\scriptsize, text=asher] at (7.6,0.35) {merge};
\end{tikzpicture}

\vspace{12pt}

\begin{itemize}
\item A \textbf{branch} is an independent line of commits
\item Useful for trying changes without breaking \texttt{main}
\item \textbf{Merge} brings the work back together
\end{itemize}

\end{frame}
% ----------------------------------------------------
\note{Branches let you work on something without disturbing the main line. Classic use case: you want to try a new model specification but are not sure if it will pan out. Create a branch (\texttt{git checkout -b new-spec}), commit freely, and if it works you merge it back into main; if not, you just delete the branch. For solo researchers, branches are optional --- many people work only on main. For collaboration, branches are essential: each person works on their own branch and merges via pull requests on GitHub. Mention that this is a preview of more advanced Git --- we are not going to practice branching in this course, but they should know it exists.}

% ----------------------------------------------------
\begin{frame}
\frametitle{The \texttt{.gitignore} file}
\centering

\vspace{5pt}

\begin{columns}[T]
\begin{column}{0.48\textwidth}
{\footnotesize \textbf{\textcolor{accent}{Track (source files):}}}

{\footnotesize
\begin{itemize}
\item R scripts (\texttt{.R})
\item \LaTeX{} source (\texttt{.tex})
\item Makefiles
\item Documentation (\texttt{.md})
\item Small data files (\texttt{.csv})
\end{itemize}
}
\end{column}
\begin{column}{0.48\textwidth}
{\footnotesize \textbf{\red{Do NOT track:}}}

{\footnotesize
\begin{itemize}
\item Generated output (\texttt{.pdf}, \texttt{.png})
\item \LaTeX{} auxiliary files (\texttt{.aux}, \texttt{.log})
\item Large data files
\item System files (\texttt{.DS\_Store})
\item Sensitive data
\end{itemize}
}
\end{column}
\end{columns}

\vspace{12pt}

{\footnotesize
\begin{itemize}
\item[] \texttt{.gitignore:}
\item[] \texttt{*.pdf}
\item[] \texttt{*.aux}
\item[] \texttt{*.log}
\item[] \texttt{.DS\_Store}
\item[] \texttt{output/}
\end{itemize}
}

\end{frame}
% ----------------------------------------------------
\note{The .gitignore file tells Git which files to ignore. The principle: track source files, not generated files. If a file can be reproduced by running a script, it does not need to be in Git. This keeps the repository small and clean. LaTeX auxiliary files (.aux, .log, .synctex.gz) are generated every time you compile and should always be ignored. Output folders (tables, figures) are generated by R scripts and should also be ignored. Large data files should not be in Git because Git is not designed for binary files and the repository will grow very quickly. System files like .DS\_Store (Mac) and Thumbs.db (Windows) should always be ignored. Show students the .gitignore from the workflow\_example project as a minimal example.}

% ----------------------------------------------------
\begin{frame}
\frametitle{Good commit habits}
\centering

\begin{itemize}[<+->]
\item Commit \textbf{often}, in small logical units
  \begin{itemize}
  \item Not: one commit per day with everything
  \item Yes: one commit per completed task or fix
  \end{itemize}
\vspace{8pt}
\item Write \textbf{meaningful commit messages}:
  \begin{itemize}
  \item \red{Bad}: \texttt{"update"}, \texttt{"changes"}, \texttt{"asdf"}
  \item \textcolor{accent}{Good}: \texttt{"Add robust SE to main models"}
  \item \textcolor{accent}{Good}: \texttt{"Fix merge bug in cleaning script"}
  \end{itemize}
\vspace{8pt}
\item Your commit history is a \textbf{lab notebook}
  \begin{itemize}
  \item You should be able to reconstruct what you did and why
  \end{itemize}
\end{itemize}

\end{frame}
% ----------------------------------------------------
\note{The commit message is the most underappreciated part of Git. Good messages make the history useful: you can look back and see when you added a feature, when you fixed a bug, and what the reasoning was. Bad messages (``update'', ``stuff'', ``final'') make the history useless. The rule of thumb: the message should complete the sentence ``This commit will...'' So ``Add robust SE to main models'' reads as ``This commit will add robust SE to main models.'' Committing often in small units means each commit is easy to understand and easy to revert if something goes wrong. One giant commit that changes 15 files is hard to review and hard to undo.}

% ----------------------------------------------------
\begin{frame}
\frametitle{The replication package as your north star}
\centering

\begin{itemize}[<+->]
\item Journals require \textbf{replication packages} (APSR, AJPS, JOP, \ldots)
\vspace{8pt}
\item From \textbf{day one} of a project, imagine you will submit one
  \begin{itemize}
  \item What would a reviewer need to go from raw data to published tables?
  \item Build that throughout --- don't assemble it in a panic at the end
  \end{itemize}
\vspace{8pt}
\item A well-organized project is already a replication package:
  \begin{itemize}
  \item Public repo with all scripts and a README
  \item Data (if shareable) or instructions to obtain it
  \item Ideally: run one command and reproduce everything
  \end{itemize}
\vspace{8pt}
\item Example: \href{https://github.com/franvillamil/vox_military}{\texttt{github.com/franvillamil/vox\_military}}
\end{itemize}

\end{frame}
% ----------------------------------------------------
\note{This is where everything comes together. The replication package is the practical test of whether your workflow is good. If a stranger can clone your repository, run \texttt{make}, and produce your tables and figures, your project is reproducible. Journals like APSR, AJPS, and JOP now require this as a condition of publication --- it is not hypothetical. The key insight is to think about replication from day one, not as an afterthought. If you organize your project well from the start, the replication package is just your project folder on GitHub. If you wait until the paper is accepted and then try to assemble one, you will spend days reconstructing which script produced which table, tracking down missing data, and documenting steps that were never written down. Point students to published replication packages on GitHub to see how other researchers structure theirs.}

% ----------------------------------------------------
\begin{frame}
\frametitle{Collaboration with Git}
\centering

\begin{itemize}[<+->]
\item Git was designed for collaboration
\vspace{8pt}
\item Basic collaboration workflow:
  \begin{itemize}
  \item Both collaborators clone the same repo
  \item Each works on their own computer
  \item \texttt{push} and \texttt{pull} to sync
  \end{itemize}
\vspace{8pt}
\item \textbf{Merge conflicts}: when two people edit the same line
  \begin{itemize}
  \item Git asks you to resolve manually
  \item Rare if you communicate and divide tasks well
  \end{itemize}
\vspace{8pt}
\item Even for solo work: syncing between laptop and desktop
\end{itemize}

\end{frame}
% ----------------------------------------------------
\note{Briefly cover collaboration because some students will co-author papers or work on group projects. The basic model is simple: both people push to and pull from the same GitHub repository. Merge conflicts happen when two people edit the same part of the same file, which Git cannot resolve automatically. In practice, merge conflicts are rare if collaborators divide the work by files (one person works on the analysis script, another on the plots script). For solo researchers, Git is still useful for syncing between multiple computers --- much more reliable than Dropbox for code projects, because Dropbox can create conflicting copies of files. Do not go deep into branches and pull requests --- that is beyond the scope of this course.}

% ====================================================
\section{Wrap-up}
% ====================================================

% ----------------------------------------------------
\begin{frame}
\frametitle{Key resources}
\centering

\begin{itemize}
\item Kieran Healy, \textit{The Plain Person's Guide to Plain Text Social Science}:
  \begin{itemize}
  \item \href{https://plain-text.co/}{\texttt{plain-text.co}}
  \end{itemize}
\vspace{8pt}
\item MIT, \textit{The Missing Semester of Your CS Education}:
  \begin{itemize}
  \item \href{https://missing.csail.mit.edu/}{\texttt{missing.csail.mit.edu}}
  \end{itemize}
\vspace{8pt}
\item Software Carpentry lessons (Unix Shell, Git):
  \begin{itemize}
  \item \href{https://software-carpentry.org/lessons/}{\texttt{software-carpentry.org/lessons}}
  \end{itemize}
\vspace{8pt}
\item Bruno Rodrigues, \textit{Building Reproducible Analytical Pipelines with R}:
  \begin{itemize}
  \item \href{https://raps-with-r.dev/}{\texttt{raps-with-r.dev}}
  \end{itemize}
\end{itemize}

\end{frame}
% ----------------------------------------------------
\note{These are all free online resources. Healy's guide is the most directly relevant: it covers plain text, LaTeX, R, Git, and project organization from a social science perspective. The MIT Missing Semester is broader (covers the command line, editors, scripting, and more) but extremely well taught. Software Carpentry has hands-on lessons designed for researchers with no programming background. Rodrigues' book covers reproducible pipelines in R specifically, including Makefiles, Docker, and functional programming. The workflow\_example repo is a minimal but complete example they can fork and adapt for their own projects.}

% ----------------------------------------------------
\begin{frame}
\frametitle{Summary}
\centering

\begin{itemize}[<+->]
\item \textbf{Code is a scientific artifact}: treat it with rigor
\vspace{8pt}
\item \textbf{Organize}: separation of concerns, clear pipeline, raw data is sacred
\item \textbf{Automate}: use Makefiles to run the pipeline
\item \textbf{Integrate}: R output feeds directly into \LaTeX{} documents
\vspace{8pt}
\item \textbf{Code well}: DRY, constants, checks, seeds, readable style
\vspace{8pt}
\item \textbf{Use plain text}: portable, versionable, scriptable
\item \textbf{Use Git}: commit often, write good messages, share via GitHub
\end{itemize}

\end{frame}
% ----------------------------------------------------
\note{Recap the main messages. The overarching theme is: your code is part of your research, and the same rigor you apply to your theory and research design should extend to your computing. None of these practices are difficult on their own; the challenge is building them into habits. The practical test: imagine a stranger needs to reproduce your results from scratch. Could they do it from your project folder? If yes, you are doing it right.}

% ----------------------------------------------------
\begin{frame}
\frametitle{For next week}
\centering

\begin{itemize}
\item Complete Assignment 10 (computing practices)
\vspace{8pt}
\item Next session: In-class exam + course review
  \begin{itemize}
  \item Exam covers sessions 1--9
  \item Second half: review and open questions
  \end{itemize}
\end{itemize}

\end{frame}
% ----------------------------------------------------
\note{Brief logistics. Assignment 10 asks students to build a small self-contained project that integrates R output into a LaTeX document. It applies the computing practices from today's lecture. Next week is the last session: the first half is the in-class exam covering sessions 1 through 9, and the second half is a course review with open questions. Remind students that the final essay is due May 20.}

% ----------------------------------------------------
\begin{frame}
\frametitle{}
\centering

Questions?

\end{frame}
% ----------------------------------------------------
\note{}
